% Capitoli V (Risultati) e VI (Discussione) — bozza in italiano, 2026-10-26, aggiornata il 2026-10-29.
% Numeri da docs/RISULTATI.md (campagna v3, tre macchine). Le celle ancora in corso
% sono marcate con \pending{...} (macro definita in main.tex).
% Da tradurre in inglese dopo la verifica.

\section{Risultati}
\label{sec:results}

Riportiamo la perdita di ricostruzione sull'holdout del modello globale rilasciato
($\mathrm{MSE}_{\mathrm{holdout}}$) e, per gli attacchi, l'AUC con il suo intervallo sui seed.
Per i canary usiamo la somma $\mathrm{AUC}_A+\mathrm{AUC}_B$ sui due bracci con ruoli
invertiti (valore atteso 1 in assenza di segnale) e, per il canary persona, l'AUC appaiato
della differenza di perdita dello stesso utente fra il braccio in cui è membro e quello in cui
non lo è. I confronti di utilità usano un test di permutazione a due code sulla media delle
perdite (10 seed per cella); gli attacchi sul canary persona un test di Mann--Whitney
unilaterale, con correzione di Holm sulle metriche secondarie.

\subsection{RQ1: rilevabilità in condizioni naturali}
\label{subsec:res_naturale}

In nessuna configurazione un attacco realistico distingue membri e non membri. Su ACN-Data,
a 10 seed, l'attacco a soglia di Yeom ha AUC 0,501 senza DP (intervallo 0,496--0,506) e 0,502
con DP locale a $\varepsilon=64$ (0,495--0,508); LiRA e RMIA con dati ausiliari del server
restano fra 0,499 e 0,501, il classificatore shadow fra 0,496 e 0,505. Su ChargePlace Scotland,
su tutte le celle (senza DP, DP per client, per sessione e per soggetto), gli stessi attacchi
stanno fra 0,486 e 0,513; lo stesso accade su ACN ridotto alle quattro feature di Scotland e
lungo la curva di capacità del modello (fino a 128--64 neuroni), dove l'AUC resta fra 0,497 e
0,508. Anche l'attacco al soggetto, che aggrega le sessioni di un utente o di un punto di
ricarica, non trova segnale: AUC 0,513 per il gruppo di 14 sessioni su ACN senza DP (1 seed su 5
significativo) e 0,520 per l'entità su Scotland (nessun seed significativo), circa 0,50 con
ciascun meccanismo di DP.

\paragraph{Avversario esterno}
Un partecipante o un osservatore esterno vede solo la sequenza dei modelli globali. Yeom e il
classificatore shadow usano già solo il modello globale, quindi i valori sopra valgono anche per
lui. Per LiRA e RMIA le ombre prendono il posto del client nell'aggregazione e gli aggiornamenti
degli altri client sono emulati con i dati ausiliari del server. Su ACN-Data (5 seed) tutte le
forme stanno fra 0,497 e 0,508 senza DP, con DP locale e con DP centrale (massimo su seed e round
0,518), come per il server sulle stesse celle: chi osserva solo il modello globale non vede più
dell'aggregatore. \pending{Scotland e controllo positivo dell'avversario esterno in corso}

\subsection{RQ1: controlli positivi}
\label{subsec:res_canary}

\paragraph{Canary per sessione}
La Tabella~\ref{tab:canary_sessione} riporta il canary per sessione (30 copie di 20 sessioni
di Office~1, 5 seed) letto sull'aggiornamento del client, cioè dalla posizione
dell'aggregatore. Senza DP il canary è riconosciuto in tutti i seed dalla perdita grezza
(somma media 1,127), da LiRA offline (1,284) e da RMIA offline (1,287). Con la DP centrale
l'aggregatore riceve l'aggiornamento soltanto clippato, perché il rumore è aggiunto dopo,
sull'aggregato: il canary è riconosciuto come senza DP (1,099, 1,280 e 1,240, tutti i seed).
Con la DP locale, dove il client aggiunge il rumore prima dell'invio, tutte le somme tornano
fra 0,999 e 1,014. Il solo clipping, quindi, non protegge dall'aggregatore; lo protegge il
rumore aggiunto dal client. Lo stesso accade su ChargePlace Scotland, con 20 sessioni di Shetland
Islands ripetute 30 volte (5 seed): al round 3 LiRA offline dà somma 1,183 senza DP e 1,145 con DP
centrale, RMIA offline 1,178 e 1,160, sempre in tutti i seed, mentre con la DP locale le somme
stanno fra 0,97 e 1,01.

\begin{table}[t]
\caption{Canary per sessione (Office~1, 30 copie, 5 seed): somma $\mathrm{AUC}_A+\mathrm{AUC}_B$
sull'aggiornamento del client (posizione dell'aggregatore), media e numero di seed con somma $>1$.
La DP centrale e quella locale usano $\varepsilon=64$ per round.}
\label{tab:canary_sessione}
\centering\footnotesize
\setlength{\tabcolsep}{4pt}
\begin{tabular}{@{}lccc@{}}
\toprule
Attacco & Senza DP & DP centrale (A2) & DP locale (A3) \\
\midrule
Perdita grezza & 1,127 (5/5) & 1,099 (5/5) & 1,014 \\
LiRA offline   & 1,284 (5/5) & 1,280 (5/5) & 0,999 \\
RMIA offline   & 1,287 (5/5) & 1,240 (5/5) & 1,012 \\
\bottomrule
\end{tabular}
\end{table}

\paragraph{Canary persona}
Il canary persona (40 utenti di JPL, 20 sessioni per utente con 30 copie, spostamento 0,5,
10 seed; Tabella~\ref{tab:canary_persona}) conferma il quadro sull'unità utente. Al round 3
l'AUC appaiato sull'aggregatore è 0,802 senza DP e 0,843 con DP centrale, significativo in
10 seed su 10, e 0,543 con DP locale (1 seed su 10, compatibile con il caso). Sul modello
rilasciato il canary è riconosciuto senza DP (0,697, 8 seed su 10) e quasi mai con DP centrale
(0,574, 2 su 10) o locale (0,555, 2 su 10): il rumore sull'aggregato protegge il modello
rilasciato, non l'aggiornamento visto dal server. Il modello rilasciato è anche ciò che vede
l'avversario esterno: il canary persona è quindi un controllo positivo anche per lui. Al round 1 senza DP l'AUC sull'aggregatore
è 0,999: è una saturazione (gli AUC dei singoli bracci sono 0,955--0,988 e il $p$ è al minimo
possibile per 20 utenti contro 20). Senza spostamento il segnale senza DP è forte al round 1
(0,775, 7 seed su 10) e si indebolisce al round 3 (0,59, 2 su 10); con la DP, e rumore
indipendente fra i bracci, non c'è in nessun regime (DP centrale 0,518 sull'aggregatore e 0,465
sul modello rilasciato, DP locale 0,484 e 0,460, nessun seed significativo su 5).
Ripetuto con rumore indipendente fra i bracci (5 seed), il quadro non cambia: DP centrale 0,711
sull'aggregatore (4 seed su 5) e 0,547 sul modello rilasciato (nessuno), DP locale 0,487 e 0,564.

% La Tabella tab:canary_persona e' quella gia' preparata (10 seed per lo spostamento 0,5).

\begin{figure*}[t]
\centering
\includegraphics[width=\textwidth]{figures/fig_canary.pdf}
\caption{Controlli positivi. A sinistra: canary per sessione, somma $\mathrm{AUC}_A+\mathrm{AUC}_B$
sull'aggiornamento del client (5 seed); la linea a 1 è il caso. A destra: canary persona, AUC
appaiato al round 3 sull'aggregatore e sul modello rilasciato (10 seed), con il numero di seed
significativi. La DP centrale lascia riconoscere il canary all'aggregatore, la DP locale no.}
\label{fig:canary}
\end{figure*}

\subsection{RQ2: costo della protezione e ottimizzazione federata}
\label{subsec:res_costo}

La Tabella~\ref{tab:utilita} riassume il costo. La DP per client a $\varepsilon=64$ per round
($\varepsilon_{\mathrm{totale}}\approx 430$ su 10 round) moltiplica la perdita sull'holdout
per 66 su ACN-Data con FedProx $\mu=0{,}01$, per 422 su ACN a quattro feature e per 536 su
Scotland. Quasi tutto il costo viene dal rumore: con il solo rumore la perdita è da 10 a 100
volte quella con il solo clipping. Lungo la griglia di $\varepsilon$ su ACN la perdita resta
fra 0,37 e 0,47 fino a $\varepsilon=16$ e scende a 0,028 (21 volte il no-DP) solo a
$\varepsilon=256$.

Senza DP FedAvg è sempre migliore di FedProx (da 2,3 a 87 volte meno perdita), e una replica
su 10 seed nuovi lo conferma in tutti i seed e in tutti i confronti. Con la DP l'ordine si
inverte. Su ACN-Data FedProx con $\mu=0{,}01$ ha 0,55 volte la perdita di FedAvg a 10 round
($p=0{,}02$); con altri 10 seed il rapporto è 0,75, e sui 20 seed insieme 0,65 ($p=0{,}008$). L'ablazione indica il meccanismo: con il solo clipping FedProx riduce la perdita
a 0,27 volte ($p<0{,}001$), con il solo rumore non cambia niente (0,85, $p=0{,}62$). Le norme degli aggiornamenti
mostrano il meccanismo: con FedAvg e la DP la norma dell'aggiornamento di ogni client resta
intorno a 8--9 dal secondo round e il clipping a $C=1$ ne lascia passare circa un sesto; con
FedProx a $\mu=0{,}01$ la norma scende a 0,6--0,7, il clipping agisce solo nell'11--16\% dei
casi e lascia passare il 97--98\% dell'aggiornamento; con $\mu=0{,}1$ non agisce più. Il
rumore, invece, ha la stessa ampiezza qualunque sia la norma, ed è per questo che FedProx non
aiuta contro il rumore. Su Scotland e su ACN a quattro feature $\mu=0{,}01$ non basta: a 10
round, sui 20 seed, il rapporto è 0,94 e 0,85 (non significativo), mentre $\mu=0{,}1$ dimezza
la perdita (0,45 e 0,53, $p<0{,}005$) e su ACN equivale a $\mu=0{,}01$ (0,66); a $\mu=0{,}1$ il
guadagno compare solo con clipping e rumore insieme, non con l'uno o l'altro da soli.
$\mu=0{,}001$ non ha effetto in nessun dataset. A 30 round il vantaggio di FedProx non si
conserva: 0,80 su ACN ($p=0{,}46$) e 0,74--0,77 su Scotland ($p>0{,}3$).
L'andamento per round spiega perché. Senza DP entrambi gli algoritmi convergono e FedAvg
scende più in fretta e più in basso. Con la DP nessuno dei due converge: FedAvg non scende mai
sotto la perdita del primo round, mentre FedProx raggiunge un minimo fra il terzo e il quinto
round (mediana 0,04--0,07, da due a quattro volte sotto FedAvg allo stesso round) e poi risale,
perché il rumore accumulato prevale; a 30 round le due curve tornano vicine
(Figura~\ref{fig:convergenza}). A un orizzonte di 5 round FedProx è migliore di FedAvg in tutti e
tre i dataset e con entrambi i valori di $\mu$: rapporto 0,59 e 0,57 su ACN-Data ($p=0{,}012$ e
$0{,}006$), 0,39 e 0,29 su Scotland ($p=0{,}004$ e $0{,}0007$), 0,56 e 0,31 su ACN a quattro
feature ($p=0{,}04$ e $0{,}0015$), per $\mu=0{,}01$ e $\mu=0{,}1$; tutti i sei confronti restano
significativi con la correzione di Holm. Poiché l'orizzonte era stato fissato dopo aver visto le
curve, abbiamo ripetuto il confronto su 10 seed nuovi con ipotesi, test e seed dichiarati prima
(Tabella~\ref{tab:replica5}): FedProx ha una perdita fra 0,30 e 0,46 volte quella di FedAvg in
tutti e sei i confronti ($p\le0{,}003$ con Holm), migliore di FedAvg in 9 o 10 seed su 10. Con la
DP conviene quindi un numero di round piccolo, che riduce anche l'$\varepsilon$ totale. La
Figura~\ref{fig:ablazione} riassume l'ablazione.

\begin{table}[t]
\caption{Replica confermativa al round 5 (10 seed nuovi, DP locale $\varepsilon=64$ per round,
$C=1$): perdita di FedProx divisa per quella di FedAvg; fra parentesi il $p$ di permutazione a
una coda corretto con Holm sui sei confronti e i seed in cui FedProx è migliore.}
\label{tab:replica5}
\centering\footnotesize
\setlength{\tabcolsep}{4pt}
\begin{tabular}{@{}lcc@{}}
\toprule
Dataset & $\mu=0{,}01$ & $\mu=0{,}1$ \\
\midrule
ACN-Data      & 0,40 (0,0005; 9/10)  & 0,36 (0,0004; 10/10) \\
Scotland      & 0,46 (0,003; 9/10)   & 0,32 (0,0005; 10/10) \\
ACN 4 feature & 0,36 (0,0005; 10/10) & 0,30 (0,0004; 10/10) \\
\bottomrule
\end{tabular}
\end{table}

\paragraph{Ruolo delle feature}
Togliendo da ACN una feature alla volta (10 seed per variante) il quadro non cambia. Senza DP
FedProx aggiunge una perdita quasi costante in valore assoluto (circa $7\cdot10^{-4}$ con
$\mu=0{,}01$), e il rapporto con FedAvg è più alto solo quando il compito è più facile: senza
i kWh richiesti, i minuti disponibili o l'ora la perdita di FedAvg scende di 3,4--3,8 volte e il
rapporto sale da 2,3 a 5--6. Con la DP il guadagno di FedProx (0,54--0,75 con $\mu=0{,}01$,
0,51--0,63 con $\mu=0{,}1$) e le norme degli aggiornamenti (7,9--8,9 con FedAvg, 0,64--0,74 con
FedProx) sono simili qualunque feature si tolga.

\begin{figure*}[t]
\centering
\includegraphics[width=\textwidth]{figures/fig_convergenza.pdf}
\caption{Perdita sull'holdout del modello rilasciato per round (mediana su 10 seed), con DP per
client a $\varepsilon=64$ per round e, tratteggiata in grigio, senza DP. Su ACN-Data FedProx con
$\mu=0{,}1$ ha solo 10 round. Senza DP la perdita scende fino al round 30; con la DP FedProx
raggiunge un minimo fra il terzo e il quinto round e poi risale.}
\label{fig:convergenza}
\end{figure*}

\begin{figure}[t]
\centering
\includegraphics[width=\columnwidth]{figures/fig_ablazione.pdf}
\caption{Rapporto fra la perdita di FedProx ($\mu=0{,}01$) e quella di FedAvg a 10 round, con
solo clipping, solo rumore e DP completa (10 seed). Simboli pieni: $p<0{,}05$; vuoti: differenza
non significativa. FedProx aiuta contro il clipping, non contro il rumore.}
\label{fig:ablazione}
\end{figure}

\begin{table*}[t]
\caption{Perdita sull'holdout del modello rilasciato (media su 10 seed, 10 round, $C=1$,
$\varepsilon=64$ per round). Fra parentesi il rapporto con FedAvg ($\mu=0$) e il $p$ del test
di permutazione.}
\label{tab:utilita}
\centering\footnotesize
\setlength{\tabcolsep}{4pt}
\renewcommand{\arraystretch}{1.15}
\begin{tabular}{@{}llccc@{}}
\toprule
Dataset & Variante & FedAvg ($\mu=0$) & FedProx $\mu=0{,}01$ & FedProx $\mu=0{,}1$ \\
\midrule
ACN-Data & Senza DP & $5{,}1\cdot10^{-4}$ & $1{,}2\cdot10^{-3}$ (2,33; $<$0,001) & -- \\
 & Solo clipping & $1{,}15\cdot10^{-2}$ & $3{,}2\cdot10^{-3}$ (0,27; $<$0,001) & -- \\
 & Solo rumore & 0,110 & 0,093 (0,85; 0,62) & -- \\
 & DP completa & 0,142 & 0,078 (0,55; 0,02) & 0,089 (0,63; 0,08) \\
\midrule
Scotland & Senza DP & $6{,}6\cdot10^{-5}$ & $3{,}0\cdot10^{-4}$ (4,5; $<$0,001) & $1{,}3\cdot10^{-3}$ (20; $<$0,001) \\
 & Solo clipping & $5{,}4\cdot10^{-3}$ & $2{,}4\cdot10^{-3}$ (0,45; 0,006) & $8{,}9\cdot10^{-3}$ (1,66; 0,01) \\
 & Solo rumore & 0,106 & 0,196 (1,85; 0,05) & 0,164 (1,55; 0,23) \\
 & DP completa & 0,183 & 0,159 (0,87; 0,56) & 0,090 (0,49; 0,015) \\
\midrule
ACN 4 feature & Senza DP & $2{,}0\cdot10^{-5}$ & $3{,}2\cdot10^{-4}$ (16; $<$0,001) & $1{,}7\cdot10^{-3}$ (87; $<$0,001) \\
 & Solo clipping & $4{,}7\cdot10^{-3}$ & $1{,}55\cdot10^{-3}$ (0,33; 0,03) & $8{,}5\cdot10^{-3}$ (1,81; 0,10) \\
 & Solo rumore & 0,162 & 0,172 (1,06; 0,86) & 0,156 (0,96; 0,91) \\
 & DP completa & 0,209 & 0,137 (0,65; 0,06) & 0,098 (0,47; 0,003) \\
\bottomrule
\end{tabular}
\end{table*}

\subsection{RQ3: unità protetta e superfici di osservazione}
\label{subsec:res_unita}

Il costo cresce con l'ampiezza dell'unità protetta. Su ACN, a 10 seed, la DP per sessione costa da 2,7 a
3,9 volte il no-DP per $\sigma$ da 1 a 5, la DP per soggetto da 9 a 11 volte, la DP per
client 62 volte nella forma centrale e 66 in quella locale; su Scotland la DP per sessione con
$\sigma=1$ costa 8 volte, quella per soggetto con $\sigma=5$ 72 volte, quella per client 536
volte. La DP per sessione costa poco ma non
protegge l'utente frequente: con la privacy di gruppo, a $\sigma=5$ ($\varepsilon\approx1{,}2$
per Caltech e JPL) il $\delta$ di gruppo supera $10^{-2}$ già con 6 sessioni e supera 1 con 14,
la mediana delle sessioni per utente in ACN-Data. La superficie conta quanto il meccanismo:
la DP centrale protegge il modello rilasciato ma non l'aggiornamento osservato dal server
(Sezione~\ref{subsec:res_canary}); la DP locale protegge entrambi. Con la DP per soggetto il canary persona non è
riconosciuto: al round 3 l'AUC appaiato sull'aggregatore è 0,499 con $\sigma=5$ ($\varepsilon\approx29$
per soggetto), 0,522 con $\sigma=8$ e 0,476 con $\sigma=1000$, senza seed significativi (5 seed), contro
0,80 senza DP nella stessa configurazione. Con la DP per sessione allo stesso $\sigma=5$
($\varepsilon\le1{,}3$ per sessione) il canary persona resta visibile all'aggregatore: AUC
appaiato 0,740 al round 1 (4 seed su 5) e 0,590 al round 3 (1 su 5), mentre con la DP per
soggetto è 0,516 e 0,499 (nessun seed). È la verifica sperimentale dell'argomento sulla privacy
di gruppo: la garanzia per sessione non copre un utente con molte sessioni.

\section{Discussione}
\label{sec:discussion}

\paragraph{Un AUC vicino a 0,5 non certifica la privacy}
In condizioni naturali nessun attacco vede i membri, con o senza DP. Senza i controlli positivi
questo risultato sarebbe ambiguo: potrebbe voler dire che il modello non memorizza o che gli
attacchi non sono in grado di estrarre il segnale. I canary mostrano che il sistema di audit
vede l'appartenenza quando la memorizzazione è indotta, e che la vede attraverso le stesse
superfici e gli stessi attacchi usati nel regime naturale. L'assenza di segnale va quindi letta
come assenza di memorizzazione misurabile nel regime operativo, non come garanzia: la garanzia
viene solo dal meccanismo, dalla sua unità protetta e dalla superficie che il rumore copre.

\paragraph{La posizione del rumore decide chi è protetto}
La DP centrale e la DP locale hanno lo stesso $\varepsilon$ ma proteggono avversari diversi.
Con la DP centrale il server riceve l'aggiornamento clippato e il canary è riconosciuto come
senza DP: il clipping limita la sensibilità ma non nasconde il contributo. Un operatore che
non si fida dell'aggregatore deve aggiungere il rumore sul client, con un costo di utilità
più alto.

\paragraph{FedProx compensa il clipping, non il rumore}
Il vantaggio di FedProx con la DP nasce dall'interazione con il clipping: il termine
prossimale riduce la norma degli aggiornamenti e quindi la parte tagliata. Poiché la maggior
parte del costo viene dal rumore, FedProx riduce il costo al più della metà e non lo avvicina
al no-DP. $\mu=0{,}1$ è uguale o migliore di $\mu=0{,}01$ su tutti e tre i dataset, mentre
$\mu=0{,}01$ basta solo su ACN; senza DP FedProx peggiora sempre: la scelta dell'algoritmo
dipende dalla presenza della DP. Il vantaggio è netto nei primi round (confermato a 5 round da
una replica dichiarata prima), si riduce a 10 e non è più significativo a 30, perché il rumore
accumulato domina.

\paragraph{Due insidie dell'audit empirico}
Nel canary persona i due bracci usavano lo stesso seed e quindi, finché il rumore della DP
veniva dal generatore globale, lo stesso rumore: la differenza appaiata lo cancellava e
mostrava un segnale anche con $\varepsilon\approx0{,}08$. Un disegno appaiato deve tenere
comuni l'inizializzazione e i dati ma usare rumore indipendente fra i bracci, come accade fra
due addestramenti reali. Inoltre, sui canary, LiRA e RMIA nella versione online danno somme
sotto 1 (fra 0,77 e 0,84). Il server dispone di pochi dati propri (112 sessioni per Office~1),
quindi l'universo delle ombre online è fatto per circa due terzi di copie dei canary: le ombre
che includono un canary lo memorizzano molto più del modello bersaglio, dove il canary è
diluito fra i dati del sito e mediato con gli altri siti. La perdita del bersaglio sul canary
cade allora dalla parte delle ombre OUT e il rapporto di verosimiglianza si inverte; le
varianti che usano solo ombre che non hanno visto il candidato non hanno questo problema.
Una diagnostica sul seed 42, che salva per i canary la statistica del bersaglio e le medie delle
ombre, mostra due effetti: le ombre che includono un canary lo memorizzano quanto o più del
bersaglio, e il bersaglio, addestrato su tutti i dati del sito, ricostruisce anche i non membri
molto meglio delle ombre, che il server addestra su pochi dati generici. Rispetto alle ombre IN i
non membri sembrano quindi più membri dei membri, ed è questo confronto che pesa nel denominatore
di RMIA online. Con le ombre addestrate sull'universo del client, a parità di seed, l'inversione
sparisce (RMIA online 1,27--1,41). Negli attacchi realistici con pochi dati del server, quindi,
le forme online vanno usate con cautela quando i candidati sono record ripetuti.

\paragraph{La DP per sessione non basta per gli utenti frequenti}
Il canary persona mostra in pratica quello che la privacy di gruppo dice in teoria: con
$\varepsilon\le1{,}3$ per sessione l'aggregatore riconosce ancora un utente con molte sessioni,
mentre la DP per soggetto allo stesso livello di rumore lo nasconde. In dati di ricarica, dove
l'utente mediano ha 14 sessioni, l'unità da proteggere è l'utente.

\paragraph{Limiti}
Le federazioni hanno tre partecipanti per dataset; il modello è un autoencoder piccolo; i
canary usano una configurazione indotta (7 feature, encoder 32--16, 3 round da 1000 epoche)
diversa da quella principale; un $\varepsilon$ di 64 per round (430 in totale) è un livello di
rumore, non una garanzia significativa; il canary persona è costruito solo su JPL; il
canary per sessione, mediato con gli altri siti, sul modello globale si vede appena anche senza
DP (somme 1,00--1,09 con Yeom e shadow) \pending{LiRA e RMIA dell'avversario esterno sul canary per sessione in corso}.
